\chapter{Игроку}
\label{ch:player}

\section{Как попасть в игру}
\label{sec:player:join}

Есть два пути, и выбирает его мастер.

\paragraph{Приглашение в лобби}
Мастер присылает ссылку вида \texttt{/lobby/<id>}. Открываете, попадаете в комнату ожидания --- это самый частый случай.

\paragraph{Заявка на персонажа}
Если мастер набирает игроков открыто, подайте заявку на \texttt{/applications/new}: опишите, кем хотите играть. Заявка попадёт в очередь к мастеру, он одобрит или отклонит её. При одобрении для вас создаётся персонаж. Свои заявки и их статус видно на \texttt{/applications}.

\section{Лобби}
\label{sec:player:lobby}

В лобби собираются перед началом. Здесь видно, кто уже подключился, и здесь мастер раздаёт персонажей.

\screenshotOptional{user/lobby.png}{комната ожидания перед игрой}

Пока мастер не назначил вам персонажа, делать в лобби нечего --- это нормально, ждите. Когда мастер запустит игру, вас перебросит в сессию автоматически; страницу перезагружать не нужно.

\section{Экран игры}
\label{sec:player:session}

Адрес --- \texttt{/session/<id>}. Экран показывает то, что мастер открыл на текущий момент:

\begin{itemize}
\item текущую сцену и её описание;
\item локацию, а если у неё есть карта --- карту с открытыми областями;
\item лист вашего персонажа;
\item объект, который мастер сейчас показывает всем;
\item ленту событий: чьи-то броски, действия, сообщения.
\end{itemize}

\screenshotOptional{user/session-player.png}{экран игрока во время сессии}

Всё обновляется само, в реальном времени. Если картинка ``замерла'' --- скорее всего, оборвалась связь; см.\ раздел \ref{sec:trouble:connection}.

\section{Действия и броски}
\label{sec:player:actions}

Действие --- не мгновенная команда, а процедура из нескольких шагов. Пока она идёт, её видят все участники, и каждый на своём шаге может что-то добавить. Ниже --- как это устроено в Dungeon World, на примере хода (\emph{move}).

\subsection{Шаг 1. Объявить ход}
Выбираете ход из доступных вашему персонажу и указываете цель --- на кого или на что он направлен. После этого действие становится видно остальным.

\subsection{Шаг 2. Бонусы и помощь}
До броска можно добавить модификаторы. Другие игроки на этом же шаге могут помочь --- их вклад войдёт в общий результат.

\subsection{Шаг 3. Бросок}
Кости бросает система. Результат попадает в ленту событий и в вашу историю бросков.

\subsection{Шаг 4. Исход}
По сумме броска определяется результат: полный успех, частичный успех или провал. Дальше выбирается конкретное последствие --- иногда автоматически, иногда выбор делает мастер или вы.

\subsection{Шаг 5. Применение}
Система собирает список изменений (урон, потраченные ресурсы, изменения счётчиков) и применяет их к персонажам и миру. При необходимости бросается урон.

\subsection{Отмена действия}
\label{sec:player:cancel}

Начатое действие можно отменить, \textbf{пока не сделан бросок}. Кнопка отмены доступна тому, кто действие начал.

После броска отменить нельзя --- кнопка просто исчезает. Это не сбой: результат уже выпал и записан в историю, отменять его задним числом было бы нечестно по отношению к остальным. Если отменить всё же нужно --- например, ход был объявлен по ошибке, --- попросите мастера: у него такая возможность остаётся.

\section{Сообщения от мастера}
\label{sec:player:messages}

Мастер может прислать сообщение лично вам --- его не увидят остальные. На такое сообщение можно ответить, и ответ тоже будет виден только мастеру. Так решаются ситуации вроде ``ты один заметил движение в кустах''.

\section{История бросков}
\label{sec:player:rolls}

Страница \texttt{/me/rolls} хранит все ваши броски: что бросали, что выпало, в какой игре. История сохраняется и после окончания подхода.
